\section{MIDI}\label{midi}

Helicase LXR has two MIDI interfaces: standard DIN connectors and a
class-compliant USB port. Both provide an input and output. Messages the
LXR generates are sent on \textbf{both} ports.

\subsection{Settings}\label{settings}

\begin{paramtable}
snc & Sync source & \texttt{off usb din pls aut}. MIDI realtime is used for \texttt{usb}, \texttt{din} and \texttt{aut}. In \texttt{aut} the priority from low to high is internal, USB, DIN, jack pulse. \\
ch  & Global channel (G) & \texttt{off} or 1 to 16. With \texttt{off}, nothing is received or sent on a global channel. \\
txf & Send filter & What the LXR sends from itself: any combination of N (notes), R (realtime), C (control) and P (program and bank change). \\
rxf & Receive filter & What the LXR acts on, with the same letters. \\
mrt & Routing & Pass-through routing: \texttt{off UtD DtD DtU bth DtA AtD all}. \\
\end{paramtable}

\attention{A card without these settings loads with receive, send and routing all \texttt{off}. The
LXR will not respond to, send or forward any MIDI until you set \texttt{rxf} and \texttt{txf} to
\texttt{all}.}

\textbf{Filter letters.} N is notes. R is realtime (clock and transport,
plus song position and MTC). C is control: NRPN, learned CCs, CC1 and
CC2. P is program change and bank change (CC0). There is no separate
bank change enable.

\textbf{Routing (\texttt{mrt}).} \texttt{UtD} USB to DIN; \texttt{DtD}
DIN to DIN; \texttt{DtU} DIN to USB; \texttt{bth} both of USB to DIN and
DIN to USB; \texttt{DtA} DIN to DIN and DIN to USB; \texttt{AtD} USB to
DIN and DIN to DIN; \texttt{all} all three. USB to USB is never routed.
A message the LXR acts on is \textbf{not} forwarded; everything else is.
Clock and transport messages are always forwarded according to
\texttt{mrt} whether or not they were used.

\subsection{Channels}\label{channels}

\begin{itemize}
\tightlist
\item
  \textbf{G} is the global channel from \texttt{ch}.
\item
  Each track has its own channel on the STEP page, per Scene:
  \texttt{off}, \texttt{glb} (follow G) or 1 to 16. A track with
  \texttt{ch\ off} has no MIDI channel at all.
\item
  A \textbf{voice channel} is a channel that some track uses that is not
  G. Tracks 6 and 7 both play voice 6 and move together.
\end{itemize}

\subsection{Notes}\label{notes}

\subsubsection{The note map}\label{the-note-map}

Each track can have a \textbf{root note} (\texttt{nte} on the STEP
page): the lowest note of its receive range. The root plays the voice's
default pitch, and higher notes in the range transpose up by the
difference.

\begin{itemize}
\tightlist
\item
  Tracks with a root own the notes from their root up to one below the
  next track's root (or 127).
\item
  Tracks without a root share what is left of the channel. On \textbf{G}
  the notes are spread in a \textbf{stripe}: note \(n\) triggers one of
  the rootless tracks in rotation at the default pitch. On a
  \textbf{voice channel} the notes are split into contiguous
  \textbf{blocks}, one per rootless track, each playing chromatically
  from its block start.
\item
  A note triggers at most one track. The pitch is clamped to 64
  semitones above the root.
\item
  \texttt{nte\ off} removes the track from the note map but keeps its
  channel for CC and NRPN.
\end{itemize}

\figplaceholder{note_map_example.png}{A worked note map.}{Keyboard strip from note 0 to 127 on channel 10 (G) with tracks 3 and 5 having roots 48 and 60. Show notes 0 to 47 striped across tracks 1, 2, 4, 6 and 7, notes 48 to 59 playing track 3 chromatically, and notes 60 to 127 playing track 5.}

\subsubsection{What a note does}\label{what-a-note-does}

\begin{itemize}
\tightlist
\item
  A note-on is a trigger of the resolved track at the note's velocity
  and pitch. A muted track does not sound, though the note is still
  consumed.
\item
  A note-off ends a held roll, but a single hit ignores it.
\item
  With a BAR overlay on, a note performs the overlay action (flam or
  held roll) instead, at the note's own pitch, in any mode.
\item
  On the Alkaline voice, velocity 101 to 127 is an accent.
\item
  With record on, notes are recorded with their true arrival time.
\end{itemize}

\subsection{Control changes}\label{control-changes}

\subsubsection{Reserved CCs}\label{reserved-ccs}

\begin{paramtable}
0 & Bank change & CC0 value $v$ on G loads Bank slot $v$. \\
1 & Morph & On G the global Morph (never recorded). On a voice channel that voice's Morph (recorded). \\
2 & BAR overlay & On G only: 0 off, 1 the \texttt{<} overlay, 2 the \texttt{>} overlay. \\
6, 38 & Data entry & NRPN value. CC38 is accepted and ignored. \\
32 & Bank select LSB & Ignored. \\
96, 97 & Data increment, decrement & NRPN step $\pm1$. \\
98, 99 & NRPN number & LSB and MSB. \\
100, 101 & RPN number & Deselects NRPN. \\
120 to 127 & Channel mode & Not used for parameters. \\
\end{paramtable}

All other CCs (3 to 5, 7 to 31, 33 to 95, 102 to 119) can be assigned
with \textbf{MIDI Learn}.

\subsubsection{NRPN}\label{nrpn}

The \textbf{NRPN number is the parameter's canonical 9-bit ID}: send
CC99 as ID shifted right by seven and CC98 as ID and 127, then CC6 for
the value. All values are absolute 7-bit values, clamped to the
parameter's own range. Morph amounts are the only exception: values 0 to
126 are doubled and 127 becomes 255.

% ============================================================
%  Helicase LXR MIDI Tables
%  The old LXR CC list is gone: parameters are addressed by NRPN =
%  canonical 9-bit target ID.  Maintain the ID ranges here.
%
%  Environments (from preamble.tex): miditableCC, miditableNRPN
%  Row format:   ID range & Meaning & Channel \\
% ============================================================

\subsection*{NRPN address space}

\begin{miditableNRPN}
0 -- 383   & Voice parameters: slot $\times$ 64 + parameter index (slots 0 to 5 = voices 1 to 6) & G or own voice channel \\
384 -- 389 & Voice Morph amounts \texttt{1vm} to \texttt{6vm} & G or own voice channel \\
390        & Retired (the former global sample-rate control); ignored & --- \\
391        & Track 7 decay (\texttt{7dc}), the open hi-hat decay & G or voice 6 channel \\
392 -- 397 & Voice outputs \texttt{1ou} to \texttt{6ou} & G or own voice channel \\
398 -- 403 & Voice effect sends \texttt{1fx} to \texttt{6fx} & G or own voice channel \\
404        & Effect Morph \texttt{fxm} & G only \\
405 -- 411 & Track length, tracks 1 to 7 & G or own voice channel \\
412 -- 418 & Track step scale, tracks 1 to 7 & G or own voice channel \\
419 -- 425 & Track shuffle, tracks 1 to 7 & G or own voice channel \\
426 -- 447 & Unused (ignored) & --- \\
448 -- 510 & Effect parameters: Effect-local index 0 to 62 & G only \\
511        & Off (ignored) & --- \\
\end{miditableNRPN}

\noindent\verify{The individual IDs from 391 to 403 are derived from the documented order
(voice Morph, \texttt{7dc}, outputs, sends) and the documented position of \texttt{fxm} at 404 and the
retired global control at 390. Confirm them against the project's published NRPN chart.}

\vspace{1ex}
\noindent The voice parameter chart (index within each Instrument type) is maintained separately by
the project. The Alkaline voice has 30 parameters (indices 0 to 29); indices 30 to 63 of its block
are unused. For the other types the indices follow the descriptor order of each Instrument.
\verify{Include the per-type voice NRPN chart in the final manual (published separately as the project's NRPN table).}

\subsection*{Parameter value notes}

\begin{miditableCC}
\texttt{osc2\_mod\_type} (Drum \texttt{mod}) & 0 \texttt{fm}, 1 \texttt{mix}, 2 \texttt{snc}, 3 to 127 \texttt{rng} & --- \\
\texttt{osc1\_type} (\texttt{typ})           & 0 \texttt{wav}, 1 \texttt{smp}, 2 \texttt{slp}, 3 \texttt{wtb}, 4 \texttt{pwm} & --- \\
\texttt{osc1\_run} (\texttt{run})            & 0 \texttt{rst}, 1 \texttt{fre}, 2 \texttt{itp}; above 2 acts as \texttt{itp} & --- \\
\texttt{voice\_chain} (\texttt{chn})         & 0 \texttt{T+D}, 1 \texttt{WFd}; above 1 acts as \texttt{WFd} & --- \\
Alkaline waveform                            & 0 \texttt{Saw}, 1 \texttt{Sqr}, 2 \texttt{Tri} & --- \\
Alkaline slide                               & 0 off, 1 on & --- \\
Alkaline tuning                              & 0 to 127, 63 is centre ($\pm$700 cents) & --- \\
Morph amounts                                & 0 to 126 are doubled; 127 gives 255 & --- \\
\end{miditableCC}

\subsection*{Default note map}

For a card with no per-track channel or note settings, the global channel stripes the notes across
the seven tracks at the default pitch, so notes 0 to 6 trigger tracks 1 to 7 in turn and the
pattern repeats every seven notes. Set a track's root note on the STEP page (\param{nte}) to give
it a chromatic range.

\verify{The earlier LXR responded on notes 36 to 42 by default. Confirm the Helicase default.}


\textbf{Instrument-agnostic.} A voice ID addresses whichever parameter
sits at that position in the slot's current Instrument type, so the same
NRPN means different things on different types. Unused IDs are ignored
and forwarded.

\textbf{Which channels work.} On G every ID works. On a voice channel
only that voice's own IDs work (its parameters, its voice Morph, output,
effect send and the timing of its tracks). Effect IDs and the other
Scene IDs are ignored on voice channels.

\textbf{A MIDI write is a menu edit.} It has the priority of a menu edit
against automation, is fanned out through the voice edit mask, always
writes the Normal value (never a Morph endpoint), is recorded in LATCH
or OVERDUB, and is never echoed back out.

\subsubsection{MIDI Learn}\label{midi-learn}

Learn maps any non-reserved CC on one channel kind to one parameter.
There are up to 128 entries.

\begin{enumerate}
\def\labelenumi{\arabic{enumi}.}
\tightlist
\item
  \textbf{Arm:} press SHIFT, then Play, then Rec, holding the earlier
  buttons. SHIFT, Play and Rec flash. While armed these buttons do none
  of their usual functions.
\item
  \textbf{Touch a parameter:} turn a knob or edit a cell. It shows
  \texttt{LRN} and its value is not changed. Cells with no MIDI address
  (tempo, settings, step data) edit normally and Learn keeps waiting.
\item
  \textbf{Send a CC.} The entry is stored and \texttt{midilrn.cfg} is
  written at once.
\item
  \textbf{Cancel} with SHIFT, Play or Rec, or by changing mode or page.
\end{enumerate}

A new entry when all 128 are used is refused (\texttt{FUL} for one
second). One CC can control many parameters; one parameter has at most
one CC. An entry applies only while its Instrument or Effect type is
loaded. Learn has no Morph endpoint: in a Morph view it still maps the
parameter itself. There is no remove in the UI; delete the line from
\texttt{midilrn.cfg}.

\figplaceholder{midi_learn_gesture.png}{Arming MIDI Learn.}{Three-step diagram of the SHIFT, Play and Rec buttons held in order with the flashing LEDs, then a display showing LRN, then a CC message arriving.}

\subsection{Program change and bank
change}\label{program-change-and-bank-change}

\begin{itemize}
\tightlist
\item
  \textbf{Program change 0 to 15} selects Scenes 1 to 16.

  \begin{itemize}
  \tightlist
  \item
    On \textbf{G} it is a whole Scene change exactly like a PERF SEQ
    press, with the usual division and reset rules.
  \item
    On a \textbf{voice channel} every track on that channel switches to
    that Scene, like holding the voice button and pressing SEQ.
  \end{itemize}
\item
  A program change to a Scene that is not loaded, or above 15, is not
  used and is forwarded.
\item
  \textbf{Bank change:} CC0 on G selects Bank slot 0 to 127 and loads it
  in the background; playback continues through Scene 17 (see the
  Storage chapter). A missing slot fails.
\item
  While a Bank is loading, program change and other non-essential
  messages are dropped.
\end{itemize}

\subsection{Clock and transport}\label{clock-and-transport}

{\def\LTcaptype{none} % do not increment counter
\begin{longtable}[]{@{}
  >{\raggedright\arraybackslash}p{(\linewidth - 2\tabcolsep) * \real{0.5000}}
  >{\raggedright\arraybackslash}p{(\linewidth - 2\tabcolsep) * \real{0.5000}}@{}}
\toprule\noalign{}
\begin{minipage}[b]{\linewidth}\raggedright
Message
\end{minipage} & \begin{minipage}[b]{\linewidth}\raggedright
Behaviour
\end{minipage} \\
\midrule\noalign{}
\endhead
\bottomrule\noalign{}
\endlastfoot
Clock & Drives the sequencer when \texttt{snc} selects its source (24
PPQ). \\
Start & Restarts from the Pattern start. \\
Continue & Resumes from where it stopped or where a song position
located it. \\
Stop & Stops and remembers the resume position. \\
Song position & Honoured while stopped. \\
MTC & Starts when eight zero quarter-frames arrive; stops after a 100 ms
timeout. \\
MMC & Not supported. \\
\end{longtable}
}

When slaved to an external clock the LXR does not generate its own;
rolls, shuffle and micro-timing stay even inside each clock period. When
the LXR is the clock master it sends clock and start, continue and stop.
Received start, stop and continue are forwarded per \texttt{mrt} but not
re-sent.

\subsection{What the LXR sends}\label{what-the-lxr-sends}

\begin{itemize}
\tightlist
\item
  \textbf{Notes (N):} every trigger the LXR plays, on the track's
  channel, with the note converted back through the note map so recorded
  output plays back identically. A track that cannot express the pitch
  in its own range sends nothing for it.
\item
  \textbf{Realtime (R):} clock and transport while the master.
\item
  \textbf{Control (C):} parameters you change by hand (not automation
  playback), as the learned CC if there is one, otherwise as NRPN on G.
  With G off, voice parameters go as NRPN on the voice channel. At most
  one parameter message is sent every 10 ms so notes and clock are never
  delayed.
\item
  \textbf{Program (P):} a program change on G for a whole-Scene change,
  program changes on the track channels for per-track changes, and CC0
  after a successful Bank load. Changes that arrived over MIDI are never
  echoed.
\end{itemize}

\subsection{Jacks}\label{jacks}

The CLK IN, RST IN and CLK OUT jacks are described in the Overview. With
\texttt{snc} set to \texttt{pls} or \texttt{aut} a clock pulse on CLK IN
drives the sequencer.

\verify{Clock pulse resolution on CLK IN and CLK OUT (the firmware defaults to a pre-4 PPQ setting) and any setting that changes it.}
